%
% Documento: Metodologia
%
\chapter{Metodologia}\label{chap:metodologia}

Este capítulo descreve as etapas metodológicas adotadas para caracterizar sinais de faltas de alta impedância em vegetação e avaliar a capacidade discriminativa dos atributos extraídos em um problema de classificação supervisionada. O desenvolvimento foi integralmente realizado em Python, utilizando as bibliotecas NumPy, SciPy, PyWavelets, scikit-learn e XGBoost.

\section{Base de dados e origem dos registros reais}\label{sec:met_dataset}

O contexto que motivou a criação da base de dados utilizada neste trabalho
remonta aos incêndios florestais de \textit{Black Saturday}, ocorridos em
fevereiro de 2009 no estado de Victoria, Austrália, que resultaram em 173
mortes e na destruição de mais de 2.000 residências
\cite{vic_bushfire_report}. Investigações posteriores indicaram que falhas em
ativos elétricos, incluindo contatos entre condutores de distribuição e
vegetação, figuraram entre as possíveis causas de alguns dos focos de incêndio.

A base de dados empregada neste trabalho é o VeHIF (\textit{Vegetation High-Impedance Fault dataset}). Os registros foram obtidos a partir do \textit{Vegetation Ignition Testing Program} (VITP), conduzido em uma rede de distribuição real de 22~kV, trifásica a três fios, no estado de Victoria, Austrália \cite{8, Gomes2021VeHIF}. Cada registro armazena sinais de corrente e tensão capturados simultaneamente por dois canais independentes:

\begin{itemize}
    \item \textbf{Canal LF (\textit{low frequency})}: registra a componente fundamental da forma de onda de corrente, sendo adequado para análise de harmônicas, inter-harmônicas e estatísticas da forma de onda no entorno dos 60~Hz.
    \item \textbf{Canal HF (\textit{high frequency})}: registra as componentes espectrais de alta frequência associadas ao arco elétrico, com foco na faixa de 2 a 10~kHz identificada na literatura como assinatura característica de FAIs em vegetação \cite{1, 8}.
\end{itemize}

O conjunto disponibilizado por \citeonline{Gomes2021VeHIF} no formato HDF5 (\textit{Hierarchical Data Format version 5}) reúne mais de 900 registros, cobrindo múltiplas espécies vegetais nativas australianas, diferentes níveis de umidade e configurações variadas de contato. O acesso ao \textit{dataset} é realizado via biblioteca \texttt{h5py} em Python. Cada arquivo contém os \textit{arrays} de sinal dos canais LF e HF, a frequência de amostragem correspondente e metadados do experimento (espécie vegetal, condição de umidade e rótulo de evento).

A inspeção inicial da estrutura dos arquivos HDF5 constitui a primeira etapa prática do trabalho e é necessária para confirmar os nomes dos \textit{datasets} internos, as taxas de amostragem efetivas e os metadados disponíveis antes de implementar o módulo de carregamento de dados.

\section{Definição do problema e das classes}\label{sec:met_classes}

O problema é formulado como uma tarefa de \textbf{classificação supervisionada binária}. O objetivo é treinar modelos capazes de distinguir, a partir de atributos extraídos dos sinais de corrente, registros de contato de vegetação com risco de ignição daqueles sem risco ou correspondentes à operação normal da rede.

A definição das classes baseia-se no critério de risco de ignição estabelecido na literatura \cite{Ozansoy2020, 8}. Conforme descrito em \citeonline{Ozansoy2023}, o processo de ignição em contatos de vegetação evolui em estágios sequenciais — contato inicial, expulsão de umidade e carbonização —, cada um com assinaturas elétricas distintas nos sinais de corrente. Registros nos quais o contato evoluiu até o estágio de carbonização ou produziu correntes de arco de amplitude e duração suficientes para causar ignição são classificados como \textbf{perigosos}; os demais, incluindo contatos de baixa energia e registros de operação sem falta, são classificados como \textbf{seguros}.

A divisão concreta das classes será determinada durante a análise exploratória, com base na estrutura de rótulos presente nos metadados do VeHIF e nos critérios operacionais descritos em \citeonline{Ozansoy2020}. Caso a distribuição entre classes seja desequilibrada, esse aspecto será considerado tanto na escolha das métricas de avaliação (Seção~\ref{sec:met_metricas}) quanto nos parâmetros de ponderação dos classificadores.

\section{Pré-processamento e preparação dos sinais}\label{sec:met_preprocessamento}

Antes da extração de atributos, cada sinal de corrente é submetido a uma cadeia de pré-processamento com o objetivo de remover artefatos e garantir a comparabilidade entre registros de condições distintas. As etapas aplicadas são as seguintes:

\begin{enumerate}
    \item \textbf{Remoção de componente DC}: subtração da média do sinal para eliminar o \textit{offset} de corrente contínua residual.

    \item \textbf{Filtragem passa-faixa}: aplicação de um filtro Butterworth de fase zero para isolar a banda de frequência de interesse em cada canal. Para o canal LF, a banda é centrada na componente fundamental de 60~Hz; para o canal HF, a faixa de 2 a 10~kHz é preservada. Os parâmetros exatos do filtro (ordem e frequências de corte) são ajustados com base na análise espectral conduzida na etapa exploratória.

    \item \textbf{Segmentação por ciclos}: o sinal filtrado é dividido em janelas temporais de comprimento fixo, correspondendo a $N_c$ ciclos da frequência fundamental. O comprimento de cada janela em amostras é dado por

    \begin{equation}
        L = N_c \cdot \left\lfloor \frac{f_s}{f_0} \right\rfloor,
        \label{eq:window_length}
    \end{equation}

    onde $f_s$ é a frequência de amostragem e $f_0 = 60$~Hz. O valor de $N_c$ é escolhido em função da duração típica dos registros e do compromisso entre resolução temporal e robustez estatística dos atributos.

    \item \textbf{Normalização}: cada janela é normalizada pelo valor de pico (intervalo $[-1,\,1]$) para atenuar diferenças de amplitude entre registros de diferentes condições de contato.
\end{enumerate}

Registros com valores ausentes, saturação de sinal ou duração inferior a dois ciclos são identificados e descartados durante a etapa de análise exploratória. Todo o pré-processamento é implementado no módulo \texttt{src/preprocessing.py}.

\section{Extração de atributos e caracterização das FAIs em vegetação}\label{sec:met_features}

Os atributos são extraídos de cada janela de sinal nos domínios do tempo, da frequência e tempo-frequência. A motivação para cada grupo de atributos está fundamentada nas características elétricas das FAIs em vegetação discutidas no Capítulo~\ref{chap:revbibliografica}.

\subsection{Atributos no domínio do tempo}

Os atributos temporais capturam características estatísticas e morfológicas da forma de onda, exploradas desde os primeiros trabalhos de detecção de FAIs \cite{4, 5}:

\begin{itemize}
    \item \textbf{Energia}: soma dos quadrados das amostras da janela,
    \begin{equation}
        E = \sum_{n=0}^{L-1} x^2[n].
        \label{eq:energy}
    \end{equation}

    \item \textbf{Assimetria (\textit{skewness})}: medida de assimetria da distribuição de amplitude,
    \begin{equation}
        \gamma_1 = \frac{1}{L} \sum_{n=0}^{L-1} \left(\frac{x[n] - \mu}{\sigma}\right)^3,
        \label{eq:skewness}
    \end{equation}
    onde $\mu$ e $\sigma$ são, respectivamente, a média e o desvio padrão da janela. Valores elevados de $\gamma_1$ refletem a condução assimétrica do arco entre semiciclos positivo e negativo, característica marcante das FAIs \cite{4}.

    \item \textbf{Curtose (\textit{kurtosis})}: medida de impulsividade da distribuição,
    \begin{equation}
        \kappa = \frac{1}{L} \sum_{n=0}^{L-1} \left(\frac{x[n] - \mu}{\sigma}\right)^4 - 3.
        \label{eq:kurtosis}
    \end{equation}
    Valores elevados indicam transientes impulsivos característicos de re-ignições do arco elétrico.

    \item \textbf{Buildup}: razão entre a energia média dos ciclos subsequentes e a energia do primeiro ciclo da janela,
    \begin{equation}
        B = \frac{E_{\mathrm{resto}}}{E_{1\degree\,\mathrm{ciclo}}},
        \label{eq:buildup}
    \end{equation}
    quantificando o crescimento progressivo de corrente descrito em \citeonline{4} como uma das assinaturas mais características das FAIs em vegetação.

    \item \textbf{Rugosidade}: variância da derivada discreta do sinal,
    \begin{equation}
        R = \mathrm{Var}\!\left(x[n] - x[n-1]\right),
        \label{eq:roughness}
    \end{equation}
    capturando flutuações locais rápidas associadas à instabilidade do arco.
\end{itemize}

\subsection{Atributos no domínio da frequência}

Os atributos espectrais são calculados a partir da FFT (\textit{Fast Fourier Transform}) de cada janela de sinal, conforme definido na Seção~\ref{sec:processamento}:

\begin{itemize}
    \item \textbf{Amplitudes harmônicas}: amplitude espectral de pico em torno das frequências $k \cdot f_0$, para $k = 1, 2, \ldots, 5$, totalizando cinco atributos. O conteúdo harmônico de baixa ordem é explorado desde \citeonline{1} e \citeonline{Huang1988} como assinatura de FAIs.

    \item \textbf{Conteúdo inter-harmônico}: razão entre a energia espectral fora das bandas harmônicas e a energia total do espectro,
    \begin{equation}
        C_{\mathrm{IH}} = \frac{\sum_{k \notin \mathcal{H}} |X[k]|^2}{\sum_{k} |X[k]|^2},
        \label{eq:interharmonic}
    \end{equation}
    onde $\mathcal{H}$ representa o conjunto de índices de frequência nas vizinhanças das harmônicas de $f_0$. Esse atributo captura a contribuição das inter-harmônicas geradas pela variação do comprimento do arco \cite{Macedo2015}.

    \item \textbf{Energia de alta frequência}: razão entre a energia espectral na faixa de 2 a 10~kHz e a energia total do espectro. Este atributo é especialmente relevante para o canal HF e constitui a base do método proposto por \citeonline{8} para detecção de VeHIFs.
\end{itemize}

\subsection{Atributos no domínio tempo-frequência}

Dois métodos de análise tempo-frequência são empregados, com base na revisão conduzida na Seção~\ref{sec:processamento}:

\begin{itemize}
    \item \textbf{Energia por nível DWT}: o sinal é decomposto pela Transformada Wavelet Discreta utilizando a wavelet de Daubechies de ordem~4 (\texttt{db4}) em $J = 5$ níveis de decomposição. Os coeficientes de detalhe $d_j[k]$ de cada nível $j$ capturam sub-bandas de frequência distintas, com resolução temporal inversamente proporcional ao nível \cite{Elkalashy2007, 6}. O atributo extraído é a energia relativa de detalhe em cada nível,
    \begin{equation}
        E_j^{\mathrm{DWT}} = \frac{\sum_k d_j^2[k]}{\sum_{n} x^2[n]},\quad j = 1, \ldots, J,
        \label{eq:dwt_energy}
    \end{equation}
    totalizando $J = 5$ atributos por janela.

    \item \textbf{Energia HF média via STFT}: a STFT é calculada sobre o sinal com janela de comprimento $M$ amostras, e para cada quadro temporal $m$ computa-se a fração de energia na faixa de 2 a 10~kHz. O atributo final é a média desta fração sobre todos os $T$ quadros,
    \begin{equation}
        E_{\mathrm{HF}}^{\mathrm{STFT}} = \frac{1}{T} \sum_{m=1}^{T} \frac{\displaystyle\sum_{k:\, f_k \in [2,\,10]\,\mathrm{kHz}} |Z(m,k)|^2}{\displaystyle\sum_{k} |Z(m,k)|^2},
        \label{eq:stft_hf}
    \end{equation}
    onde $Z(m,k)$ denota o espectrograma calculado pela STFT \cite{7}.
\end{itemize}

O vetor de atributos resultante da aplicação de todos os extratores a uma janela de sinal de um único canal contém $5 + 7 + 6 = 18$ atributos. Quando os dois canais LF e HF são considerados conjuntamente, o vetor completo possui 36 atributos.

\section{Seleção de atributos}\label{sec:met_selecao}

A seleção de atributos tem duplo propósito: reduzir a dimensionalidade do vetor de entrada dos classificadores e identificar quais atributos são mais informativos para o problema em questão, contribuindo para a interpretabilidade dos resultados. Dois métodos complementares são empregados:

\begin{itemize}
    \item \textbf{F-score ANOVA}: teste univariado que avalia se a distribuição de cada atributo difere significativamente entre as classes, por meio da razão entre a variância inter-classe e a variância intra-classe.

    \item \textbf{Importância por floresta aleatória}: um classificador Random Forest é treinado e a redução média de impureza de Gini associada a cada atributo é utilizada como medida de importância global \cite{9}. Esse método captura interações entre atributos, o que o torna complementar ao F-score ANOVA de natureza univariada.
\end{itemize}

Os resultados de ambas as abordagens serão comparados para identificar consensos e discrepâncias. O número $k$ de atributos selecionados será determinado pela análise do desempenho dos classificadores em função de $k$ (diagrama de cotovelo), equilibrando desempenho e parcimônia do modelo.

\section{Classificadores supervisionados}\label{sec:met_classifiers}

Quatro algoritmos de classificação supervisionada são avaliados, abrangendo diferentes paradigmas de aprendizado e graus de interpretabilidade, conforme motivado na Seção~\ref{sec:ia}:

\begin{enumerate}
    \item \textbf{Árvore de Decisão} (\textit{Decision Tree}, DT): estrutura hierárquica de partições binárias no espaço de atributos, com alta interpretabilidade. Classifica uma instância percorrendo o grafo da raiz até um nó folha, aplicando limiares sobre atributos individuais. Referência consolidada no contexto de FAIs desde \citeonline{5}.

    \item \textbf{Máquina de Vetores de Suporte} (\textit{Support Vector Machine}, SVM): classificador de margem máxima com \textit{kernel} RBF (\textit{Radial Basis Function}) para captura de separações não lineares entre classes \cite{Aljohani2020}. Os atributos são normalizados com \textit{z-score} antes da aplicação do modelo.

    \item \textbf{XGBoost} (\textit{eXtreme Gradient Boosting}): método \textit{ensemble} baseado em gradiente que combina sequencialmente árvores de decisão rasas, cada uma corrigindo os erros residuais da anterior. Atingiu 98,17\% de acurácia em dados reais do PBSP no trabalho de \citeonline{Ma2020}.

    \item \textbf{Rede Neural Artificial} (\textit{Multilayer Perceptron}, MLP): rede com duas camadas ocultas (128 e 64 neurônios, ativação ReLU), treinada com parada precoce (\textit{early stopping}) para evitar sobreajuste. Normalização dos atributos aplicada antes do treinamento.
\end{enumerate}

Todos os classificadores são implementados via \texttt{scikit-learn} e \texttt{xgboost} com hiperparâmetros padrão como ponto de partida, permitindo uma avaliação inicial comparável. Ajuste fino via busca em grade poderá ser aplicado ao melhor modelo identificado nos experimentos preliminares.

\section{Métricas de avaliação e análise de desempenho}\label{sec:met_metricas}

O desempenho dos classificadores é avaliado por meio de validação cruzada estratificada com $k = 5$ partições. A estratificação garante que a proporção entre as classes seja preservada em cada partição, o que é especialmente relevante em conjuntos com distribuição desequilibrada.

As métricas reportadas para cada dobra, e como médias com desvio padrão, são:

\begin{itemize}
    \item \textbf{Acurácia}: proporção de classificações corretas sobre o total de instâncias,
    \begin{equation}
        \mathrm{Ac} = \frac{VP + VN}{VP + VN + FP + FN},
        \label{eq:accuracy}
    \end{equation}
    onde $VP$, $VN$, $FP$ e $FN$ denotam, respectivamente, verdadeiros positivos, verdadeiros negativos, falsos positivos e falsos negativos.

    \item \textbf{Precisão}: proporção de instâncias classificadas como positivas que são efetivamente positivas,
    \begin{equation}
        \mathrm{P} = \frac{VP}{VP + FP}.
        \label{eq:precision}
    \end{equation}

    \item \textbf{Recall} (sensibilidade): proporção de instâncias positivas reais que são corretamente identificadas,
    \begin{equation}
        \mathrm{R} = \frac{VP}{VP + FN}.
        \label{eq:recall}
    \end{equation}
    No contexto de detecção de faltas, o \textit{recall} possui especial relevância operacional: faltas não detectadas (falsos negativos) correspondem a eventos de risco de ignição não identificados.

    \item \textbf{F\textsubscript{1}-score}: média harmônica entre precisão e \textit{recall},
    \begin{equation}
        F_1 = 2 \cdot \frac{\mathrm{P} \cdot \mathrm{R}}{\mathrm{P} + \mathrm{R}},
        \label{eq:f1}
    \end{equation}
    que equilibra o impacto de falsos positivos e falsos negativos em uma única métrica agregada.
\end{itemize}

Os resultados finais são apresentados em tabelas comparativas entre os quatro classificadores, considerando cada configuração de atributos testada: canal LF, canal HF, conjunto completo de 36 atributos e subconjunto de $k$ atributos selecionados. A análise visa identificar qual combinação de atributos e classificador oferece o melhor compromisso entre desempenho preditivo e interpretabilidade para o problema de detecção de VeHIFs.
